% Auto-generated from book/05_forms_topographs.md — do not edit by hand
% Regenerated by scripts/md_to_latex.py

\chapter{Chapter 5 --- Quaternionic Forms and Flux Topographs}
\label{ch:ch05_topographs}

This chapter lifts Conway's topographs of binary quadratic forms (Hatcher Chapters 4--5) to the gauged Hopf lattice. We define \textbf{flux topographs}---value landscapes on the discrete lattice whose periodic separator structures encode stability, periodicity, and classification. The resulting theory is the direct higher-dimensional, quaternionic analogue of Hatcher's visual number theory, with Magic Islands as a natural generalization of class-number phenomena and reduced forms---stated carefully as a \textbf{Model}.

\textbf{Learning goals}

\begin{enumerate}
\item Generalize binary quadratic forms and their topographs to flux functionals on the gauged Hopf lattice.  
\item Define flux topographs, value landscapes, and separator structures.  
\item Understand periodicity, reduced configurations, and Magic Islands.  
\item See how gauge symmetries act on topographs (equivariance).  
\item Prepare the visual and arithmetic foundation for classification (Chapter 6) and the \(Z\mapsto\) flywheel map (Chapter 7).
\end{enumerate}

\textbf{Figures in this chapter}

\begin{center}
\small
\setlength{\tabcolsep}{4pt}
\begin{tabularx}{\textwidth}{@{}>{\raggedright\arraybackslash}p{0.12\textwidth}>{\raggedright\arraybackslash}X>{\raggedright\arraybackslash}X@{}}
\toprule
Tag & File & Role \\
\midrule
Fig. 5.1 & \texttt{figures/\allowbreak{}fig5\_\allowbreak{}1\_\allowbreak{}flux\_\allowbreak{}topograph\_\allowbreak{}schematic.\allowbreak{}png} & Flux topograph value landscape \\
Fig. 5.2 & \texttt{figures/\allowbreak{}fig5\_\allowbreak{}2\_\allowbreak{}separator\_\allowbreak{}structure.\allowbreak{}png} & Periodic / sign-crossing separators \\
Fig. 5.3 & \texttt{figures/\allowbreak{}fig5\_\allowbreak{}3\_\allowbreak{}magic\_\allowbreak{}island.\allowbreak{}png} & Magic Island schematic + model Z-sweep \\
Fig. 5.4 & \texttt{figures/\allowbreak{}fig5\_\allowbreak{}4\_\allowbreak{}symmetry\_\allowbreak{}on\_\allowbreak{}topograph.\allowbreak{}png} & Gauge action on a topograph \\
Aux A5.1 & \texttt{figures/\allowbreak{}aux5\_\allowbreak{}1\_\allowbreak{}topograph\_\allowbreak{}evolution.\allowbreak{}png} & Evolution under repeated right phase \\
\bottomrule
\end{tabularx}
\end{center}


\textbf{Claim discipline}

\begin{center}
\small
\setlength{\tabcolsep}{4pt}
\begin{tabularx}{\textwidth}{@{}>{\raggedright\arraybackslash}X>{\raggedright\arraybackslash}X@{}}
\toprule
Claim & Type \\
\midrule
Classical topograph properties (periodicity, arithmetic progression rule, reduced forms) from Hatcher & \textbf{Theorem} (classical; cited) \\
Flux topographs, separator structures, Magic Islands, gauge equivariance of candidate constructions & \textbf{Model} (core of OP2) \\
\texttt{qga/\allowbreak{}lib/\allowbreak{}flux\_\allowbreak{}topograph.\allowbreak{}py}, \texttt{kingdom.\allowbreak{}core.\allowbreak{}flux\_\allowbreak{}flywheel} APIs & \textbf{Software fact} \\
\bottomrule
\end{tabularx}
\end{center}


\bigskip
\noindent\rule{\textwidth}{0.4pt}
\bigskip

\section{5.1 From binary quadratic forms to flux functionals}
\label{ch:ch05_topographs:5-1-from-binary-quadratic-forms-to-flux-functionals}

In Hatcher, a binary quadratic form

\[
f(x,y) = ax^2 + bxy + cy^2
\]

acquires a \textbf{topograph}---a value landscape drawn on the dual of the Farey triangulation. Adjacent regions differ by an arithmetic progression rule, and periodic separator lines appear for indefinite forms (positive non-square discriminant).

We lift this picture to the gauged Hopf lattice of Chapters 3--4:

A \textbf{flux functional} on the lattice assigns a real (or integer) value to vertices---and, in richer versions, to edges or small cycles---such that the assignment transforms naturally under gauge actions. A \textbf{flux topograph} is the visual landscape of these values together with their level sets and separator structures.

Natural candidate functionals (all \textbf{Model} choices for OP2 experiments):

\begin{center}
\small
\setlength{\tabcolsep}{4pt}
\begin{tabularx}{\textwidth}{@{}>{\raggedright\arraybackslash}X>{\raggedright\arraybackslash}X@{}}
\toprule
Name (book helper) & Geometric meaning \\
\midrule
\texttt{norm} & \(x^2+y^2+z^2\) on pure-imaginary part (quadratic form on \(\mathrm{Im}\,\mathbb{H}\)) \\
\texttt{hopf\_\allowbreak{}y1} / \texttt{hopf\_\allowbreak{}height} & Classical Hopf base coordinates (Ch. 2) \\
\texttt{phase} & Fiber phase proxy \(\mathrm{atan2}(x_4,x_3)\) \\
portal stability & \texttt{map\_\allowbreak{}z\_\allowbreak{}to\_\allowbreak{}flywheel} scores (Ch. 7 preview) \\
\bottomrule
\end{tabularx}
\end{center}


\textbf{Model note.} The precise choice of functional is part of the definition of flux topographs and is tied to \textbf{Open Problem 2}. No uniqueness claim is made in this chapter.

\bigskip
\noindent\rule{\textwidth}{0.4pt}
\bigskip

\section{5.2 Flux topographs and separator structures}
\label{ch:ch05_topographs:5-2-flux-topographs-and-separator-structures}

\subsection{Definition (working)}
\label{ch:ch05_topographs:definition-working}

A \textbf{flux topograph} is a tuple

\[
\bigl(\Lambda,\; E,\; V,\; \mathcal{S}\bigr)
\]

where \(\Lambda\) is a discrete point set in \(S^3\), \(E\) is an edge set (along-fiber and/or inter-fiber from Ch. 3), \(V:\Lambda\to\mathbb{R}\) is a flux functional, and \(\mathcal{S}\) is the separator structure extracted from \(V\) (sign crossings, critical levels).

In code:

\begin{lstlisting}[style=qga]
FluxTopograph(points, values, edges, functional)
build_flux_topograph(points, edges=..., functional=...)
detect_separators(topo, threshold=..., mode='sign'|'level')
\end{lstlisting}

\begin{figure}[htbp]
  \centering
  \includegraphics[width=0.92\textwidth,height=0.42\textheight,keepaspectratio]{fig5_1_flux_topograph_schematic.png}
  \caption{Figure 5.1 --- Flux topograph value landscape.}
  \label{fig:fig5-1-flux-topograph-schematic}
\end{figure}

\begin{quote}\small\textit{Figure 5.1.* Value landscape for the functional \texttt{hopf\_\allowbreak{}y1} on an angle-sampled lattice, shown in stereographic \(S^3\) and pushed to the base \(S^2\).}\end{quote}

\subsection{Separator structures}
\label{ch:ch05_topographs:separator-structures}

\textbf{Separator structures} are the loci where the functional changes sign or crosses a critical threshold. On a graph, they appear as edges \((i,j)\) with \(V(i)\) and \(V(j)\) on opposite sides of a threshold. Connected components of those edges are discrete analogues of Hatcher's separator curves.

In the indefinite / oscillatory case, separators can become \textbf{periodic} under the gauge group---exactly analogous to Hatcher's periodic separator lines for indefinite binary forms, which encode infinitude (Pell-type families).

\begin{figure}[htbp]
  \centering
  \includegraphics[width=0.92\textwidth,height=0.42\textheight,keepaspectratio]{fig5_2_separator_structure.png}
  \caption{Figure 5.2 --- Separator structures.}
  \label{fig:fig5-2-separator-structure}
\end{figure}

\begin{quote}\small\textit{Figure 5.2.* Dark edges: sign-crossing separators of \texttt{hopf\_\allowbreak{}height} on the candidate adjacency graph. Component count and edge sets are experimental OP2 diagnostics, not theorems.}\end{quote}

\subsection{Arithmetic progression rule (lift)}
\label{ch:ch05_topographs:arithmetic-progression-rule-lift}

Classically, values on three regions meeting at a topograph edge satisfy a linear relation (Hatcher's arithmetic progression rule). On the gauged Hopf lattice we only partially lift this: for each edge we can measure jumps \(V(j)-V(i)\) and study their regularity under gauge moves (\texttt{arithmetic\_\allowbreak{}progression\_\allowbreak{}residuals}). A full equivariant AP rule compatible with quaternion multiplication is part of \textbf{OP2}.

\bigskip
\noindent\rule{\textwidth}{0.4pt}
\bigskip

\section{5.3 Periodicity, reduced configurations, and Magic Islands}
\label{ch:ch05_topographs:5-3-periodicity-reduced-configurations-and-magic-islands}

\subsection{Periodicity under the gauge group}
\label{ch:ch05_topographs:periodicity-under-the-gauge-group}

When a flux topograph is (approximately) periodic under a composite gauge sequence \(g\) from Chapter 4---i.e. \(g^k\cdot\Phi \approx \Phi\) for some \(k\)---the orbit corresponds to a \textbf{reduced configuration}. The book helper \texttt{periodicity\_\allowbreak{}score} reports nearest-neighbor total-space return and multiset distance of values.

\subsection{Magic Islands}
\label{ch:ch05_topographs:magic-islands}

\textbf{Magic Islands} are pockets of enhanced stability in parameter space (detuning \(\delta\omega\), gauge strength, layers, or atomic number \(Z\) in the portal Model). Inside an island the topograph (or the derived stability score) shows strong order: high \texttt{stability\_\allowbreak{}score}, noble-gas locks, low variation. Outside, scores drop and dynamical simulations become more chaotic (Ch. 3 Aux A3.1).

They are the narrative generalization of finite class number and reduced forms in Hatcher's theory. \textbf{Arithmetic characterization} (class-number-like invariants predicting islands) is deferred to Open Problem 3 / Chapter 6.

\begin{figure}[htbp]
  \centering
  \includegraphics[width=0.92\textwidth,height=0.42\textheight,keepaspectratio]{fig5_3_magic_island.png}
  \caption{Figure 5.3 --- Magic Island stability region.}
  \label{fig:fig5-3-magic-island}
\end{figure}

\begin{quote}\small\textit{Figure 5.3.* Left: schematic island in parameter space. Right: Kingdom Come Model \texttt{stability\_\allowbreak{}score} vs \(Z\) (gold lines mark noble-gas \(Z\)). \textbf{Software fact} for the right panel; physical interpretation remains \textbf{Model}.}\end{quote}

\textbf{Claim type.} Visual and software properties of Magic Islands in the portal: \textbf{Model} + \textbf{Software fact}. Classical reduced forms / class number: \textbf{Theorem} (Hatcher). Correspondence between the two: \textbf{Open} (OP3).

\bigskip
\noindent\rule{\textwidth}{0.4pt}
\bigskip

\section{5.4 Gauge equivariance of flux topographs}
\label{ch:ch05_topographs:5-4-gauge-equivariance-of-flux-topographs}

Any well-defined flux topograph must be \textbf{equivariant} under the gauge group of Chapter 4: applying a left or right multiplication must map the topograph to another topograph of the same type (possibly after transporting values or recomputing a geometric functional).

This is the direct analogue of Hatcher's requirement that topographs transform properly under linear fractional transformations.

\begin{figure}[htbp]
  \centering
  \includegraphics[width=0.92\textwidth,height=0.42\textheight,keepaspectratio]{fig5_4_symmetry_on_topograph.png}
  \caption{Figure 5.4 --- Gauge symmetry acting on a flux topograph.}
  \label{fig:fig5-4-symmetry-on-topograph}
\end{figure}

\begin{quote}\small\textit{Figure 5.4.* Left multiplication by \(i\) maps a \texttt{hopf\_\allowbreak{}y1} topograph to its image. For geometric functionals recomputed after the move, values track the geometry exactly.}\end{quote}

\begin{figure}[htbp]
  \centering
  \includegraphics[width=0.92\textwidth,height=0.42\textheight,keepaspectratio]{aux5_1_topograph_evolution.png}
  \caption{Auxiliary Figure A5.1 --- Topograph evolution.}
  \label{fig:aux5-1-topograph-evolution}
\end{figure}

\begin{quote}\small\textit{Auxiliary Figure A5.1.* Repeated right phase advances: base scatter colored by \texttt{phase} functional. Separator patterns should transform predictably if OP2 axioms hold.}\end{quote}

\textbf{Open Problem 2 (core of this chapter).} Formalize the axioms of a flux topograph so that:

\begin{enumerate}
\item The structure \textbf{reduces} to Conway/Hatcher topographs under a suitable projection or restriction.  
\item Separator periodicity and an arithmetic progression rule lift naturally.  
\item \textbf{Gauge equivariance} holds for a fixed adjacency rule from Chapter 3 (OP1).  
\item Magic Islands emerge as finite or periodic reduced regions with measurable stability.
\end{enumerate}

\textbf{Sandbox:} \texttt{qga/\allowbreak{}lib/\allowbreak{}flux\_\allowbreak{}topograph.\allowbreak{}py} --- \texttt{build\_\allowbreak{}flux\_\allowbreak{}topograph}, \texttt{detect\_\allowbreak{}separators}, \texttt{periodicity\_\allowbreak{}score}, \texttt{separator\_\allowbreak{}equivariance\_\allowbreak{}score}.

\bigskip
\noindent\rule{\textwidth}{0.4pt}
\bigskip

\section{5.5 First computational labs}
\label{ch:ch05_topographs:5-5-first-computational-labs}

Helpers: \texttt{lib/\allowbreak{}flux\_\allowbreak{}topograph.\allowbreak{}py} \(\cdot\) portal: \texttt{map\_\allowbreak{}z\_\allowbreak{}to\_\allowbreak{}flywheel} \(\cdot\) \textbf{Appendix C ?C.3}.

\begin{itemize}
\item \textbf{5.A} \texttt{build\_\allowbreak{}flux\_\allowbreak{}topograph} + \texttt{detect\_\allowbreak{}separators}.
\item \textbf{5.B} \texttt{periodicity\_\allowbreak{}score} under a gauge sequence.
\item \textbf{5.C} Z-stability via \texttt{map\_\allowbreak{}z\_\allowbreak{}to\_\allowbreak{}flywheel} (no \texttt{magic\_\allowbreak{}flag} --- use \texttt{stability\_\allowbreak{}score} / \texttt{is\_\allowbreak{}noble\_\allowbreak{}gas}).
\item \textbf{5.D} \texttt{separator\_\allowbreak{}equivariance\_\allowbreak{}score} (OP2 diagnostic).
\end{itemize}

\begin{lstlisting}[style=qga,language=python]
from lib.flux_topograph import build_flux_topograph, detect_separators
from lib.hopf_lattice import sample_angle_lattice, candidate_adjacency
pts = sample_angle_lattice(3, 8, 8)
along, inter = candidate_adjacency(pts)
topo = build_flux_topograph(pts, edges=along+inter, functional="hopf_height")
print(len(detect_separators(topo)))
\end{lstlisting}

\bigskip
\noindent\rule{\textwidth}{0.4pt}
\bigskip

\section{Exercises}
\label{ch:ch05_topographs:exercises}

\textbf{5.A (hand).} State the classical arithmetic progression rule for Hatcher topographs (or look it up in TN Ch. 4) and propose a quaternionic lift for flux values on adjacent lattice elements.

\textbf{5.B (hand).} In two sentences, explain why gauge equivariance is a necessary condition for a flux topograph to be well-defined.

\textbf{5.C (code).} Complete Labs 5.A--5.B. Report separator component count and \texttt{periodicity\_\allowbreak{}score} for your functional.

\textbf{5.D (code).} Run Lab 5.C for \(Z\in\{2,10,26,79,118\}\). Which values show high \texttt{stability\_\allowbreak{}score} or \texttt{is\_\allowbreak{}noble\_\allowbreak{}gas=True}? Record the pattern.

\textbf{5.E (visual).} Generate a flux topograph and apply one gauge transformation (Fig. 5.4). Describe how separator curves / value colors transform.

\textbf{5.F (Hatcher bridge).} In Hatcher, reduced forms and class number classify quadratic forms up to equivalence. Sketch how Magic Islands might play an analogous role for flux configurations on the gauged Hopf lattice.

\textbf{5.G (Open Problem 2 teaser).} Using \texttt{build\_\allowbreak{}flux\_\allowbreak{}topograph} and \texttt{detect\_\allowbreak{}separators}, test whether separator counts survive under \texttt{candidate\_\allowbreak{}adjacency} + left/right multiplications (\texttt{separator\_\allowbreak{}equivariance\_\allowbreak{}score}). Document failures---they inform OP2 axioms.

\textbf{5.H (forward).} Why will the classification theory of Chapter 6 depend on the periodic separator structures developed here?

\textbf{5.I (software honesty).} In one paragraph, distinguish: (i) pedagogical flux functionals in \texttt{qga/\allowbreak{}lib/\allowbreak{}flux\_\allowbreak{}topograph.\allowbreak{}py}, (ii) portal \texttt{stability\_\allowbreak{}score} from \texttt{map\_\allowbreak{}z\_\allowbreak{}to\_\allowbreak{}flywheel}, (iii) classical Conway topographs. Assign claim types.

\bigskip
\noindent\rule{\textwidth}{0.4pt}
\bigskip

\section{Code and asset pointers}
\label{ch:ch05_topographs:code-and-asset-pointers}

\begin{lstlisting}[style=qga]
qga/lib/flux_topograph.py
  build_flux_topograph, detect_separators, periodicity_score,
  apply_gauge_to_topograph, separator_equivariance_score,
  arithmetic_progression_residuals, stability_landscape_z

qga/lib/hopf_lattice.py
kingdom.core.flux_flywheel.map_z_to_flywheel[_extended]
kingdom.viz.magic_island.build_magic_island_heatmap
\end{lstlisting}

\textbf{Figures:} \texttt{scripts/\allowbreak{}generate\_\allowbreak{}ch5\_\allowbreak{}figures.\allowbreak{}py} \textbf{Related portal:} Flux Flywheel tab; Magic Island heatmap; Lattice Simulator (dynamical backdrop). \textbf{Open problems:} OP2 (\texttt{notes/\allowbreak{}open\_\allowbreak{}problems.\allowbreak{}md}); depends on OP1 adjacency.

\bigskip
\noindent\rule{\textwidth}{0.4pt}
\bigskip

\section{Looking ahead}
\label{ch:ch05_topographs:looking-ahead}

Flux topographs are the visual and arithmetic heart of Part III---the direct quaternionic lift of Hatcher's topographs, still with \textbf{Model}-level axioms pending OP2. In \textbf{Chapter 6} we classify these topographs, enumerate reduced configurations, and formalize Magic Islands as higher-dimensional analogues of class numbers and reduced forms. \textbf{Chapter 7} connects classified objects to the \(Z\mapsto\) flywheel map, completing the representation-theoretic bridge from number theory to the periodic-table proxy in the Kingdom Come Model.

With flux topographs in hand, we are ready for the classification theory that turns visual landscapes into arithmetic invariants.

\bigskip
\noindent\rule{\textwidth}{0.4pt}
\bigskip


